\documentclass[11pt,letterpaper]{article}
\usepackage{../../Common/course_notes}
\usepackage{../../Common/course_macros}

\begin{document}
\setlength{\abovedisplayskip}{6pt}
\setlength{\belowdisplayskip}{6pt}
\setlength{\textfloatsep}{10pt}
\setlength{\intextsep}{8pt}
\color{blackmain}

\begin{center}
{\LARGE\textbf{Geometría y Física I: Mecánica Clásica y Formas de Cartan}}\\[0.25cm]
{\Large\textcolor{orangeacc}{\textbf{Semana 3: El Grupo de Rotaciones $SO(2)$, $SO(3)$ y Traslaciones}}}\\[0.15cm]
{\large\textbf{Rigidez, ortogonalidad, orientación, no-conmutatividad y rotaciones infinitesimales}}\\[0.3cm]
{\large\textcolor{graymuted}{Notas de Estudio --- Módulo 0: El Arsenal Indicial y Geométrico}}\\[0.2cm]
\end{center}

\vspace{0.2cm}

% =====================================================================
\manualsection[1. Rigidez y Matrices Ortogonales]{1. La Física de la Rigidez: De la Preservación de Distancias a $R^T R = I$}

¿Qué distingue físicamente al movimiento de un cuerpo rígido (como un giroscopio, una roca o una molécula indeformable) de la deformación de un fluido o un resorte elástico? 

La respuesta fundamental es la \textbf{invarianza de las distancias internas}: si dos partículas del cuerpo están separadas por un vector $\vec{v}$, después de que el cuerpo gire el nuevo vector $\vec{v}'$ debe conservar exactamente la misma longitud. Más aún, si tomamos dos vectores cualesquiera $\vec{u}$ y $\vec{v}$ adheridos al cuerpo rígido, el ángulo relativo entre ellos no puede alterarse.

\begin{intuitionbox}{Deducción de la Condición Ortogonal}
Sea $R: \R^n \to \R^n$ una transformación lineal que asigna a cada vector $\vec{v}$ su posición rotada $\vec{v}' = R\vec{v}$.
Exigimos que para cualesquiera dos vectores $\vec{u}, \vec{v} \in \R^n$ se conserve su producto interno euclidiano:
\[
\langle R\vec{u}, R\vec{v} \rangle = \langle \vec{u}, \vec{v} \rangle.
\]
Usando la definición de la transpuesta matricial en producto interno ($\langle A\vec{x}, \vec{y} \rangle = \langle \vec{x}, A^T \vec{y} \rangle$), el miembro izquierdo se reescribe como:
\[
\langle \vec{u}, R^T (R\vec{v}) \rangle = \langle \vec{u}, \vec{v} \rangle \quad \Longleftrightarrow \quad \langle \vec{u}, (R^T R - I)\vec{v} \rangle = 0.
\]
Dado que esta igualdad debe cumplirse de manera idéntica para todos los vectores $\vec{u}, \vec{v} \in \R^n$, concluimos necesariamente:
\[
\boxed{R^T R = I} \quad \Longleftrightarrow \quad R^{-1} = R^T.
\]
\end{intuitionbox}

\begin{theorembox}{Definición Formal: Grupo Ortogonal $O(n)$}
Una matriz real $R \in \mathcal{M}_{n\times n}(\R)$ se denomina \textbf{ortogonal} si satisface $R^T R = I$. El conjunto de todas estas transformaciones forma el \textbf{Grupo Ortogonal} $O(n)$:
\[
O(n) \equiv \{ R \in \mathcal{M}_{n\times n}(\R) \mid R^T R = I \}.
\]
Las columnas de una matriz ortogonal $R = [\vec{c}_1 \mid \dots \mid \vec{c}_n]$ constituyen siempre una \textbf{base ortonormal} de $\R^n$, pues la condición matricial equivale en índices a:
\[
(R^T)_{ki} R_{kj} = R_{ik} R_{jk} = \delta_{ij} \quad \Longleftrightarrow \quad \vec{c}_i \cdot \vec{c}_j = \delta_{ij}.
\]
\end{theorembox}

% =====================================================================
\manualsection[2. Rotaciones en el Plano: El Grupo SO(2)]{2. Construcción Columna a Columna de Rotaciones en 2D: $SO(2)$}

Para encontrar la forma explícita de una rotación en el plano por un ángulo $\theta$ en sentido antihorario, no necesitamos memorizar matrices: basta seguir el destino de los dos vectores de la base canónica ortonormal $\{\hat{e}_1, \hat{e}_2\}$.

\begin{center}
\begin{tikzpicture}[scale=1.1]
    % Ejes originales
    \draw[-{Stealth}, color=graymuted] (-0.4,0) -- (3.6,0) node[right, font=\footnotesize] {$x$};
    \draw[-{Stealth}, color=graymuted] (0,-0.4) -- (0,3.3) node[above, font=\footnotesize] {$y$};

    \coordinate (O) at (0,0);
    \coordinate (E1) at (2.4, 0);
    \coordinate (E2) at (0, 2.4);

    % Base rotada
    \coordinate (RE1) at ({2.4*cos(35)}, {2.4*sin(35)});
    \coordinate (RE2) at ({-2.4*sin(35)}, {2.4*cos(35)});

    % Vectores base originales
    \draw[-{Stealth[length=2.8mm]}, line width=1.3pt, color=graymuted] (O) -- (E1) node[below=2pt, font=\small\bfseries] {$\hat{e}_1$};
    \draw[-{Stealth[length=2.8mm]}, line width=1.3pt, color=graymuted] (O) -- (E2) node[left=2pt, font=\small\bfseries] {$\hat{e}_2$};

    % Vectores rotados
    \draw[-{Stealth[length=3mm]}, line width=1.8pt, color=orangeacc] (O) -- (RE1) node[above right=1pt, font=\small\bfseries] {$R\hat{e}_1$};
    \draw[-{Stealth[length=3mm]}, line width=1.8pt, color=navyaccent] (O) -- (RE2) node[above left=1pt, font=\small\bfseries] {$R\hat{e}_2$};

    % Angulo theta para e1
    \draw[color=orangeacc, thick] (1.0, 0) arc (0:35:1.0);
    \node[color=orangeacc, font=\footnotesize\bfseries] at (1.25, 0.35) {$\theta$};

    % Angulo theta para e2
    \draw[color=navyaccent, thick] (0, 1.0) arc (90:125:1.0);
    \node[color=navyaccent, font=\footnotesize\bfseries] at (-0.35, 1.25) {$\theta$};

    % Proyecciones dashed
    \draw[dashed, color=graymuted] (RE1) -- ({2.4*cos(35)}, 0) node[midway, right=1pt, font=\scriptsize] {$\sin\theta$};
    \draw[dashed, color=graymuted] (RE1) -- (0, {2.4*sin(35)});
    \node[below=2pt, font=\scriptsize, color=graymuted] at ({2.4*cos(35)}, 0) {$\cos\theta$};
\end{tikzpicture}
\end{center}

Al rotar por un ángulo $\theta$:
\begin{align*}
R(\theta)\hat{e}_1 &= \cos\theta\,\hat{e}_1 + \sin\theta\,\hat{e}_2 = \begin{pmatrix} \cos\theta \\ \sin\theta \end{pmatrix}, \\
R(\theta)\hat{e}_2 &= -\sin\theta\,\hat{e}_1 + \cos\theta\,\hat{e}_2 = \begin{pmatrix} -\sin\theta \\ \cos\theta \end{pmatrix}.
\end{align*}
Colocando estos dos vectores columna lado a lado en la matriz:
\[
\boxed{R(\theta) = \begin{pmatrix} \cos\theta & -\sin\theta \\ \sin\theta & \cos\theta \end{pmatrix}}
\]
Calculando su determinante:
\[
\det R(\theta) = \cos^2\theta - (-\sin^2\theta) = \cos^2\theta + \sin^2\theta = +1.
\]

% =====================================================================
\manualsection[3. El Espejo y la Orientación: SO(n) vs O(n)]{3. Rotaciones Propias frente a Reflexiones: ¿Por qué $\det R = +1$?}

Si tomamos el determinante a ambos lados de la condición de rigidez $R^T R = I$:
\[
\det(R^T R) = \det(R^T)\det(R) = (\det R)^2 = \det(I) = 1 \quad \implies \quad \det R = \pm 1.
\]
Esta bifurcación divide a las transformaciones ortogonales en dos clases geométricas disjuntas:

\begin{theorembox}{Definición Formal: Grupo Especial Ortogonal $SO(n)$}
\begin{itemize}
    \item \textbf{Rotaciones Propias ($\det R = +1$):} Preservan el volumen y la orientación espacial (regla de la mano derecha). Constituyen el \textbf{Grupo Especial Ortogonal}:
    \[
    SO(n) \equiv \{ R \in O(n) \mid \det R = +1 \}.
    \]
    \item \textbf{Reflexiones e Inversiones Impropias ($\det R = -1$):} Preservan distancias pero invierten la orientación (como un espejo que cambia una mano derecha por una mano izquierda). No pueden alcanzarse mediante un giro continuo desde la identidad $I$.
\end{itemize}
\end{theorembox}

\begin{center}
\begin{tikzpicture}[scale=0.95]
    % Panel Izquierdo: Rotacion (orientacion preservada)
    \begin{scope}[xshift=-4.5cm]
        \coordinate (O) at (0,0);
        \draw[-{Stealth}, color=orangeacc, line width=1.5pt] (O) -- (2.2, 0.4) node[right, font=\footnotesize\bfseries] {$\vec{u}'$};
        \draw[-{Stealth}, color=navyaccent, line width=1.5pt] (O) -- (-0.4, 2.2) node[above, font=\footnotesize\bfseries] {$\vec{v}'$};
        \draw[color=orangeacc, ->, thick] (1.0, 0.18) arc (10:75:1.0);
        \node[font=\scriptsize\bfseries] at (1.0, 1.2) {Sentido horario/antihorario};
        \node[font=\small\bfseries, align=center] at (1.0, -0.8) {(a) Rotación Propia: $\det R = +1$\\[-2pt]{\scriptsize Mano derecha conservada}};
    \end{scope}

    % Panel Derecho: Espejo (orientacion invertida)
    \begin{scope}[xshift=3.0cm]
        \draw[line width=1.5pt, color=graymuted] (0, -0.8) -- (0, 2.7) node[above, font=\footnotesize\bfseries] {Espejo};
        \coordinate (O) at (0,0);
        \draw[-{Stealth}, color=orangeacc, line width=1.5pt] (O) -- (2.0, 0.6) node[right, font=\footnotesize\bfseries] {$\vec{u}$};
        \draw[-{Stealth}, color=navyaccent, line width=1.5pt] (O) -- (0.6, 2.0) node[above right, font=\footnotesize\bfseries] {$\vec{v}$};

        \draw[-{Stealth}, color=orangeacc, dashed, line width=1.5pt] (O) -- (-2.0, 0.6) node[left, font=\footnotesize\bfseries] {$\vec{u}^*$};
        \draw[-{Stealth}, color=navyaccent, dashed, line width=1.5pt] (O) -- (-0.6, 2.0) node[above left, font=\footnotesize\bfseries] {$\vec{v}^*$};
        \node[font=\small\bfseries, align=center] at (0, -0.8) {(b) Reflexión: $\det R = -1$\\[-2pt]{\scriptsize Mano derecha se vuelve izquierda}};
    \end{scope}
\end{tikzpicture}
\end{center}

% =====================================================================
\manualsection[4. Rotaciones en 3D: SO(3) y Matrices Canónicas]{4. Rotaciones en el Espacio: $SO(3)$ Alrededor de los Ejes Coordinados}

En el espacio tridimensional $\R^3$, rotar alrededor de un eje coordinado deja dicho eje inmóvil y actúa como una rotación bidimensional $SO(2)$ en el plano perpendicular:

\begin{align*}
R_x(\psi) &= \begin{pmatrix} 1 & 0 & 0 \\ 0 & \cos\psi & -\sin\psi \\ 0 & \sin\psi & \cos\psi \end{pmatrix}, \qquad
R_y(\theta) = \begin{pmatrix} \cos\theta & 0 & \sin\theta \\ 0 & 1 & 0 \\ -\sin\theta & 0 & \cos\theta \end{pmatrix}, \\
R_z(\phi) &= \begin{pmatrix} \cos\phi & -\sin\phi & 0 \\ \sin\phi & \cos\phi & 0 \\ 0 & 0 & 1 \end{pmatrix}.
\end{align*}
Cada una satisface $R_i^T R_i = I$ y $\det(R_i) = +1$.

% =====================================================================
\manualsection[5. No-Conmutatividad de SO(3)]{5. La No-Conmutatividad en 3D: El Experimento del Libro}

Una diferencia fundamental entre dos y tres dimensiones es que mientras dos rotaciones planas siempre conmutan ($R(\theta_1)R(\theta_2) = R(\theta_2)R(\theta_1) = R(\theta_1+\theta_2)$), \textbf{las rotaciones en el espacio tridimensional no conmutan}:
\[
R_x(\pi/2) R_z(\pi/2) \neq R_z(\pi/2) R_x(\pi/2).
\]

\begin{center}
\begin{tikzpicture}[scale=0.92]
    % Camino 1: Primero Z luego X
    \begin{scope}[xshift=-4.6cm]
        \node[font=\footnotesize\bfseries, color=orangeacc] at (1.5, 3.2) {Camino 1: Primero $R_z(\pi/2)$, luego $R_x(\pi/2)$};
        \draw[fill=creambg, draw=blackmain, line width=1pt] (0,0) rectangle (2.6, 2.0);
        \node[font=\scriptsize\bfseries] at (1.3, 1.0) {Portada arriba};
        \draw[-{Stealth}, line width=1.3pt, color=orangeacc] (2.9, 1.0) -- (4.1, 1.0) node[midway, above=1pt, font=\scriptsize] {$R_z$};
        \draw[fill=orangebg, draw=orangeacc, line width=1pt] (4.4,0) rectangle (6.4, 2.6);
        \node[font=\scriptsize\bfseries, color=navyaccent] at (5.4, 1.3) {Giro en plano};
    \end{scope}

    % Desigualdad
    \node[font=\LARGE\bfseries, color=navyaccent] at (2.6, 1.0) {$\neq$};

    % Camino 2: Primero X luego Z
    \begin{scope}[xshift=3.4cm]
        \node[font=\footnotesize\bfseries, color=navyaccent] at (1.5, 3.2) {Camino 2: Primero $R_x(\pi/2)$, luego $R_z(\pi/2)$};
        \draw[fill=creambg, draw=blackmain, line width=1pt] (0,0) rectangle (2.6, 2.0);
        \node[font=\scriptsize\bfseries] at (1.3, 1.0) {Portada arriba};
        \draw[-{Stealth}, line width=1.3pt, color=navyaccent] (2.9, 1.0) -- (4.1, 1.0) node[midway, above=1pt, font=\scriptsize] {$R_x$};
        \draw[fill=navyaccent!15, draw=navyaccent, line width=1pt] (4.4,0) rectangle (7.0, 1.2);
        \node[font=\scriptsize\bfseries, color=blackmain] at (5.7, 0.6) {Lomo arriba};
    \end{scope}
\end{tikzpicture}
\end{center}

Esta falta de conmutatividad significa algebraicamente que el grupo $SO(3)$ es \textbf{no abeliano}.

% =====================================================================
\manualsection[6. Teorema de Euler y Rotaciones Infinitesimales]{6. El Teorema de Rotación de Euler y la Velocidad Angular $\vec{\omega}$}

\manualsubsection{El Teorema de Rotación de Euler: El Eje Fijo}
Todo cuerpo que rota en el espacio tridimensional mantiene al menos una línea fija que no se mueve: el \textbf{eje de rotación}.

\begin{theorembox}{Teorema de Euler}
Para toda matriz de rotación $R \in SO(3)$, existe al menos un vector no nulo $\vec{v} \neq \vec{0}$ tal que:
\[
R\vec{v} = \vec{v}.
\]
Es decir, $\lambda = 1$ es siempre un valor propio de toda matriz de rotación en 3D.
\end{theorembox}

\noindent\textit{Demostración:} Consideremos el determinante $\det(R - I)$:
\begin{align*}
\det(R - I) &= \det(R - R R^T) = \det[R(I - R^T)] = \det(R)\det[(I - R)^T] \\
&= (+1)\det(I - R) = (-1)^3 \det(R - I) = -\det(R - I).
\end{align*}
Un número real igual a su propio negativo debe ser cero: $2\det(R - I) = 0 \implies \det(R - I) = 0$. Por el teorema fundamental del álgebra lineal, el sistema homogéneo $(R - I)\vec{v} = \vec{0}$ tiene soluciones no triviales.

\manualsubsection{Rotaciones Infinitesimales: De Matrices a la Velocidad Angular}
Consideremos una rotación muy pequeña alrededor del origen que se aparta de la identidad en un parámetro infinitesimal $\delta\theta \ll 1$:
\[
R(\delta\theta) \approx I + \Omega\,\delta\theta.
\]
Si exigimos la condición de rigidez $R^T R = I$:
\[
(I + \Omega^T\,\delta\theta)(I + \Omega\,\delta\theta) = I + (\Omega + \Omega^T)\delta\theta + \mathcal{O}(\delta\theta^2) = I.
\]
Para que el término de primer orden se anule, la matriz generadora $\Omega$ debe ser \textbf{antisimétrica}:
\[
\boxed{\Omega^T = -\Omega} \quad \Longleftrightarrow \quad \Omega_{ij} = -\Omega_{ji}.
\]
Toda matriz real antisimétrica de $3 \times 3$ contiene solo 3 parámetros independientes $(\omega_1, \omega_2, \omega_3)$:
\[
\Omega = \begin{pmatrix} 0 & -\omega_3 & \omega_2 \\ \omega_3 & 0 & -\omega_1 \\ -\omega_2 & \omega_1 & 0 \end{pmatrix}.
\]
Al multiplicar por un vector de posición $\vec{r} = (x, y, z)^T$:
\[
\Omega\vec{r} = \begin{pmatrix} -\omega_3 y + \omega_2 z \\ \omega_3 x - \omega_1 z \\ -\omega_2 x + \omega_1 y \end{pmatrix} = \vec{\omega} \times \vec{r}, \qquad \text{donde } \vec{\omega} = \begin{pmatrix} \omega_1 \\ \omega_2 \\ \omega_3 \end{pmatrix}.
\]
Así, la derivada temporal de la posición de un punto en un sólido rígido en rotación pura es exactamente:
\[
\boxed{\dot{\vec{r}} = \vec{\omega} \times \vec{r}}
\]
La matriz antisimétrica $\Omega$ y el producto vectorial $\vec{\omega}\times$ representan el mismo objeto físico.

% =====================================================================
\manualsection[7. Rotaciones vs. Traslaciones]{7. Rotaciones frente a Traslaciones: Dos Naturalezas Geométricas}

Para concluir, contrastamos las rotaciones con las traslaciones espaciales:
\begin{center}
\begin{tabular}{|l|c|c|}
\hline
\textbf{Propiedad} & \textbf{Rotación $R \in SO(3)$} & \textbf{Traslación $T_{\vec{a}}$} \\ \hline
Definición & $\vec{x}' = R\vec{x}$ & $\vec{x}' = \vec{x} + \vec{a}$ \\ \hline
Linealidad & Lineal ($R(\alpha\vec{u}+\beta\vec{v})=\alpha R\vec{u}+\beta R\vec{v}$) & No lineal si $\vec{a}\neq\vec{0}$ (afín) \\ \hline
Punto fijo & Siempre fija el origen $R\vec{0} = \vec{0}$ y todo el eje $\vec{v}$ & Ningún punto fijo si $\vec{a}\neq\vec{0}$ \\ \hline
Álgebra de composición & No conmutativa en 3D ($R_1 R_2 \neq R_2 R_1$) & Siempre conmutativa ($T_{\vec{a}}T_{\vec{b}} = T_{\vec{b}}T_{\vec{a}}$) \\ \hline
\end{tabular}
\end{center}

\end{document}
